\documentclass[11pt]{article}
\usepackage[T1]{fontenc}
\usepackage{lmodern}
\usepackage[margin=1.05in]{geometry}
\usepackage{amsmath,amssymb,amsthm}
\usepackage[hidelinks]{hyperref}

% Fill these in before release. Empty fields produce no byline.
\newcommand*{\AuthorName}{}
\newcommand*{\AuthorEmail}{}
\newcommand{\byline}{\AuthorName\ifx\AuthorEmail\empty\else
  \ifx\AuthorName\empty\else\\\fi
  {\small\href{mailto:\AuthorEmail}{\expandafter\nolinkurl\expandafter{\AuthorEmail}}}\fi}
\newtheorem{theorem}{Theorem}[section]
\newtheorem{lemma}[theorem]{Lemma}
\newtheorem{corollary}[theorem]{Corollary}
\title{Large prime divisors of Fermat congruences\\and semiprime pseudoprimes}
\author{\byline}
\date{}

\begin{document}
\maketitle
\begin{abstract}
We count the odd integers \(n\le X\) with exactly two prime factors,
counted with multiplicity, that satisfy \(2^{n-1}\equiv1\pmod n\).
Their number is \(O(X^{3/4}/(\log X)^{5/4})\). For distinct
factors \(r,s\), we group the pairs by \(\gcd(r-1,s-1)\).
A bound on the number of prime divisors of a single Mersenne number
handles small gcds, and two applications of the Brun--Titchmarsh
inequality handle large gcds. We also bound the number of composite
\(n\le X\) having a prime \(q>Y\) that divides both \(n\) and
\(2^{n-1}-1\) by \(O((X/Y)^2/\log Y)\). A large prime factor
of \(n-1\) gives a further bound for prime powers common to
\(n\) and \(a^{n-1}-1\) at any fixed base \(a\ge2\). These estimates quantify
the composite contribution to two possible counts of twin primes.
A lower bound large enough to force infinitely many twin primes remains open.
\end{abstract}

\section{Semiprime pseudoprimes}

An odd composite integer \(n\) is a base-\(2\) Fermat pseudoprime if
\(n\mid 2^{n-1}-1\). Write \(Q_2(X)\) for the number of such integers
\(n\le X\) that are products of two primes, allowing equal factors.
All logarithms are natural except where the base is indicated. The
constants in \(O(\cdot)\) and \(\ll\) are absolute unless a
dependence is indicated.

\begin{theorem}\label{thm:semiprime}
For \(X\ge3\),
\[
             Q_2(X)\ll \frac{X^{3/4}}{(\log X)^{5/4}}.
\]
\end{theorem}

Luca, Ralaivaosaona and Jim\'enez Urroz proved
\(Q_2(X)\ll X^{4/5}/(\log X)^{2/5}\) in
\cite[Theorem~3.1]{LRJ}. Their result concerns products up to \(X\)
with unrestricted factor sizes. Their separate estimate for factors
in a common dyadic interval uses a different counting parameter.
Theorem~\ref{thm:semiprime} sharpens their bound for the unrestricted
count. We use their parametrization by the common divisor of the two
shifted prime factors and count the small values of that parameter
one at a time.

We use \(\varphi\) for Euler's totient function, \(\mu\) for the
M\"obius function and \(w(h)=h/\varphi(h)\).
The following two estimates keep track of both primality conditions
in the large-gcd range.

\begin{lemma}\label{lem:BT}
For all positive integers \(h\) and real \(T\ge2\),
\[
 \#\{1\le t\le T: ht+1\text{ is prime}\}
       \ll w(h)\frac{T}{\log(2T)}.
\]
For positive integers \(a\) and real \(M\ge2\),
\[
 \sum_{\substack{M<g\le2M\\ ag+1\ \mathrm{prime}}}w(g)
        \ll w(a)\frac{M}{\log(2M)}.
\]
In addition \(\sum_{a\le Z} w(a)/a\ll\log(2Z)\) for \(Z\ge1\).
\end{lemma}

\begin{proof}
The first inequality is the Brun--Titchmarsh inequality applied
to the class \(1\bmod h\) up to \(hT+1\) \cite{MV}.
For the second, use
\[
            w(g)=\sum_{d\mid g}\frac{\mu^2(d)}{\varphi(d)}.
\]
If \(d\le\sqrt M\), write \(g=dt\) and apply the first estimate
with \(h=ad\) and \(T=2M/d\). Since \(w(ad)\le w(a)w(d)\),
the total contribution of these \(d\) is at most a constant times
\[
 \frac{M w(a)}{\log(2M)}
      \sum_{d\ge1}\frac{\mu^2(d)}{\varphi(d)^2}.
\]
The series converges by its Euler product. For \(d>\sqrt M\),
count all multiples of \(d\) up to \(2M\). The inequality
\(\varphi(d)\ge\sqrt{d/2}\), obtained prime power by prime power,
bounds their contribution by
\(2M\sum_{d>\sqrt M}1/(d\varphi(d))\ll M^{3/4}\).
This gives the second assertion. Finally, expanding \(w(a)\) in the
same way and summing the harmonic series gives
\[
  \sum_{a\le Z}\frac{w(a)}{a}
       \le (1+\log Z)\sum_{d\ge1}\frac{\mu^2(d)}{d\varphi(d)}
       \ll\log(2Z).
\]
\end{proof}

\begin{proof}[Proof of Theorem~\ref{thm:semiprime}]
Prime squares contribute at most \(\sqrt X\). For the remaining
integers write \(n=rs\), with \(r<s\) odd primes, and set
\(g=\gcd(r-1,s-1)\). The order of \(2\) modulo \(r\) divides
both \(r-1\) and \(rs-1\), whose gcd is \(g\). Exchanging \(r\)
and \(s\) gives
\[
                  r,s\mid 2^g-1,\qquad r,s\equiv1\pmod g.
\]
For fixed \(g\ge2\), let \(t_g\) count the prime divisors of
\(2^g-1\) congruent to \(1\bmod g\). Their product is at least
\((g+1)^{t_g}\) and at most \(2^g-1\), whence
\(t_g\ll g/\log(2g)\). The pairs with \(g\le G\) therefore
contribute
\begin{equation}\label{eq:small}
     \sum_{2\le g\le G}t_g^2\ll \frac{G^3}{(\log(2G))^2}.
\end{equation}
For example, splitting the sum at \(\sqrt G\) gives the bound directly.

Take \(G=(X\log X)^{1/4}\) and assume \(X\) is large enough
that \(G<X^{3/8}\). For \(G<g\le X^{3/8}\), divide the range into
intervals \(M<g\le2M\), starting at \(M=G\) and truncating the
last one. Put \(r=ag+1\), \(s=bg+1\). Then
\[
       a\le \sqrt X/g,\qquad b\le B:=X/(ag^2),\qquad
       B\ge \sqrt X/g\ge X^{1/8}.
\]
For each eligible \(g,a\), Lemma~\ref{lem:BT} bounds the choices
of \(b\) for which \(bg+1\) is prime by
\(\ll Xw(g)/(ag^2\log X)\). Summing first over \(g\), while
retaining the requirement that \(ag+1\) be prime, bounds the number
of pairs in this block by
\begin{align*}
   &\ll \frac{X}{M^2\log X}
       \sum_{a\le\sqrt X/M}\frac1a
       \sum_{\substack{M<g\le2M\\ag+1\ \mathrm{prime}}}w(g)\\
   &\ll \frac{X}{M\log X\log(2M)}
               \sum_{a\le\sqrt X/M}\frac{w(a)}a
     \ll \frac{X}{M\log X}.
\end{align*}
Here \(\log(2M)\gg\log X\) because \(M\ge G\). The sum over
these blocks is \(O(X/(G\log X))\).

For \(g>X^{3/8}\), count every pair of positive integers
\(a,b\) with \(ab\le X/g^2\). There are at most
\((X/g^2)(1+\log X)\) such pairs. Summing over \(g\) gives
\(O(X^{5/8}\log X)\). At our choice of \(G\), this last
term and the prime squares are absorbed by the claimed bound,
and both \eqref{eq:small} and \(X/(G\log X)\) have the
required order.
\end{proof}

\section{One large common prime divisor}

A single prime factor of \(n\) can satisfy the Fermat congruence
even when \(n\) fails it. The following estimate applies to all
composites, with any number of prime factors.

\begin{theorem}\label{thm:common}
Let \(2\le Y\le X\). The number of pairs \((n,q)\) with
\(n\le X\) composite, \(q>Y\) prime and
\(q\mid\gcd(n,2^{n-1}-1)\) is
\[
                        \ll \frac{(X/Y)^2}{\log Y}.
\]
If \(n\ge5\) is odd, then \(n\) is prime if and only if
\(\gcd(n,2^{n-1}-1)\) has a prime factor exceeding
\(2n/\log_2 n\).
\end{theorem}

\begin{proof}
Write \(n=qk\). Since \(n-1=qk-1\equiv k-1\pmod{q-1}\),
Fermat's little theorem gives
\begin{equation}\label{eq:cofactor}
             q\mid 2^{n-1}-1\quad\Longleftrightarrow\quad
             q\mid 2^{k-1}-1.
\end{equation}
If \(n\) is composite, \(2\le k\le X/Y\). Distinct primes
\(q>Y\) dividing \(2^{k-1}-1\) have product at most \(2^{k-1}-1\).
Their number is at most \((k-1)\log 2/\log Y\). Summing over
\(k\) counts pairs and proves the first assertion.

If \(q>2n/\log_2 n\), then \(k<\tfrac12\log_2 n\).
Also \(2n/\log_2 n>\sqrt n\) for \(n\ge5\), as follows
from \(\log_2 n<2\sqrt n\).
For \(k\ge2\) this gives \(0<2^{k-1}-1<\sqrt n<q\), contrary
to \eqref{eq:cofactor}. Thus \(k=1\). Conversely, if \(n\ge5\)
is prime, \(q=n\) divides the gcd and exceeds the indicated
threshold.
\end{proof}

\section{A large prime divisor of the exponent}

The prime \(q\) in Theorem~\ref{thm:common} divides the modulus
\(n\). We now examine an independent condition on the exponent
\(n-1\).

\begin{lemma}\label{lem:exponent}
Let \(n\ge3\) and \(n-1=k\ell\), where \(\ell>\sqrt n\) is
prime. Then \(n\) is prime if and only if \(\ell\mid\varphi(n)\).
Equivalently, some residue \(u\not\equiv1\pmod n\) satisfies
\(u^\ell\equiv1\pmod n\). If \(n\) is composite, every integer
\(a\) satisfies
\[
    \gcd(n,a^{n-1}-1)=\gcd(n,a^k-1).
\]
\end{lemma}

\begin{proof}
If \(n\) is prime, its cyclic unit group has order \(k\ell\) and
contains an element of order \(\ell\). Conversely, any nontrivial
\(\ell\)-th root has order \(\ell\) in the group of units modulo
\(n\), so \(\ell\mid\varphi(n)\). Since \(\ell\nmid n\), the
product formula for \(\varphi(n)\) then gives a prime \(r\mid n\)
with \(r\equiv1\pmod\ell\). The cofactor \(n/r\) is also \(1\bmod
\ell\) and satisfies \(1\le n/r<\ell\). Thus \(n=r\).

Now let \(n\) be composite and let \(r^j\mid n\) be any prime
power. If \(r\mid a\), neither power of \(a\) in the displayed
gcd identity is congruent to \(1\bmod r\). Otherwise the order
of \(a\bmod r^j\) divides \(\varphi(r^j)\), which divides
\(\varphi(n)\). That order is coprime to \(\ell\), so it divides
\(k\ell\) exactly when it divides \(k\). This proves equality
of all prime-power divisibility conditions in the two gcds.
\end{proof}

\begin{theorem}\label{thm:joint}
Fix an integer \(a\ge2\), \(\theta>1/2\) and \(\delta>1-\theta\).
The number of composite \(n\in(X,2X]\) for which \(n-1\) has a
prime factor \(\ell>X^\theta\) and \(\gcd(n,a^{n-1}-1)\) has a
prime-power divisor at least \(X^\delta\) is
\[
  \ll_{a,\theta,\delta}
       \frac{X^{2-\theta-\delta}+X^{2-2\theta}}{\log X}.
\]
\end{theorem}

\begin{proof}
If \(\theta>1\), no such \(\ell\) exists for large \(X\). We may
therefore take \(1/2<\theta\le1\), so \(\delta>0\).
For sufficiently large \(X\), we have \(\ell>\sqrt{2X}\).
Write \(n=k\ell+1\), so \(k\le2X^{1-\theta}\). By
Lemma~\ref{lem:exponent}, the prime power in the hypothesis divides
\(a^k-1\). For each prime \(r\) for which some power at least
\(X^\delta\) divides \(a^k-1\), take the least such power \(s_r\).
The \(s_r\) are pairwise coprime, so their product is at most
\(a^k-1\). Their number is at most \(k\log a/(\delta\log X)\).

At least one of these \(s_r\) divides \(n\). For fixed \(k,s_r\),
the congruence \(k\ell\equiv-1\pmod{s_r}\) is impossible if
\(k\) and \(s_r\) have a common prime divisor. Otherwise it specifies
one residue class for \(\ell\). The condition \(X<k\ell+1\le2X\)
allows at most \(X/(k s_r)+1\le X^{1-\delta}/k+1\) integers
in that class. Summing this bound over the selected prime powers
and over \(k\le2X^{1-\theta}\) proves the theorem.
\end{proof}

\section{Consequences for prime pairs}

Averaging over the base gives a different measure of the contribution
from semiprimes. For odd \(n\ge3\), let
\[
   F(n)=\frac1n\#\{a\bmod n:a^{n-1}\equiv1\pmod n\}.
\]
Every residue satisfying the congruence is a unit. The denominator
counts all residue classes, including nonunits.

\begin{lemma}\label{lem:average}
The sum of \(F(n)\) over odd products of two primes \(n\le X\),
counting each integer once and allowing squares, is \(O(\sqrt X)\).
\end{lemma}

\begin{proof}
For \(n=rs\) with distinct odd primes, cyclicity of the unit groups
modulo \(r\) and \(s\) and the Chinese remainder theorem give
\(F(rs)=g^2/(rs)\), where \(g=\gcd(r-1,s-1)\).
Write \(r=ag+1\) and \(s=bg+1\) with \(a<b\). Then
\(F(rs)\le1/(ab)\) and \(g\le\sqrt{X/(ab)}\). Summing
over all positive \(a,b\) gives an upper bound of
\[
   \sqrt X\sum_{a,b\ge1}(ab)^{-3/2}\ll\sqrt X.
\]
Prime squares contribute at most \(\sqrt X\), since \(F(n)\le1\).
\end{proof}

Let \(T(X)\) count odd primes \(n\le X\) for which \(n-2\) is
prime. Let \(B(X)\) count odd integers \(n\le X\) having at most
two prime factors counted with multiplicity, with \(n-2\) prime
and \(n\mid 2^{n-1}-1\). Theorem~\ref{thm:semiprime} gives
\[
             0\le B(X)-T(X)\ll X^{3/4}(\log X)^{-5/4}.
\]
Likewise, let \(H_{\theta}(X)\) count odd \(n\le X\) with
\(n-2\) prime and a prime divisor \(q>X^\theta\) of
\(\gcd(n,2^{n-1}-1)\), where \(1/2<\theta<1\) is fixed.
For \(X^\theta\ge2\), Theorem~\ref{thm:common} gives
\[
 0\le H_{\theta}(X)-\bigl(T(X)-T(X^\theta)\bigr)
          \ll \frac{X^{2(1-\theta)}}{\theta\log X}.
\]
Every prime \(n>X^\theta\) counted by \(T(X)\) supplies the divisor
\(q=n\). Every other integer contributing to \(H_\theta(X)\)
is composite. Both upper bounds discard the condition that \(n-2\)
is prime when estimating this composite contribution.

For the corresponding average over bases, Lemma~\ref{lem:average}
and \(F(n)=1-1/n\) at an odd prime give
\[
  \sum_{\substack{3\le n\le X,\ n\ \mathrm{odd}\\
                n-2\ \mathrm{prime},\ \Omega(n)\le2}} F(n)
       =T(X)+O(\sqrt X),
\]
where \(\Omega(n)\) counts prime factors with multiplicity.
The discarded \(1/n\) terms contribute at most \(1+\log X\).

At the logarithmic cutoff of Theorem~\ref{thm:common}, one qualifying
integer \(n=p+2\) for arbitrarily large primes \(p\) would suffice
for infinitude. The divisor identity counts the exceptions when a
smaller cutoff is used. Proving the existence of qualifying divisors
at the logarithmic cutoff remains the missing step.

\section*{Acknowledgments}
AI assistants contributed to drafting this manuscript and checking
its mathematical arguments.

\begin{thebibliography}{9}
\bibitem{LRJ}
F.~Luca, D.~Ralaivaosaona and J.~Jim\'enez Urroz,
\emph{On pseudoprime RSA moduli}, Afrika Matematika
\textbf{36} (2025), article 118.
\href{https://doi.org/10.1007/s13370-025-01338-1}{doi:10.1007/s13370-025-01338-1}.

\bibitem{MV}
H.~L.~Montgomery and R.~C.~Vaughan,
\emph{The large sieve}, Mathematika \textbf{20} (1973), 119--134.
\href{https://doi.org/10.1112/S0025579300004708}{doi:10.1112/S0025579300004708}.
\end{thebibliography}
\end{document}
